% Candidate mathematics; no empirical validity is asserted by these identities.
\section*{Protected interaction under a shared physical texture}
Fix a simulator state, instruction, proprioception, and policy inference noise.
Let $A_P(u,v)\in\mathbb R^d$ be the selected action coordinates from rendering
path $P\in\{R,N\}$, when the two views use texture parameters $u$ and $v$.
The physical deployment is the diagonal $u=v=\theta$. Off-diagonal evaluations
are camera-input interventions, not additional realizable physical textures.
Define
\[
 d_{1,P}=A_P(\theta,0)-A_P(0,0),\quad
 d_{2,P}=A_P(0,\theta)-A_P(0,0),\quad
 d_{J,P}=A_P(\theta,\theta)-A_P(0,0),
\]
\[
 I_P=A_P(\theta,\theta)-A_P(\theta,0)-A_P(0,\theta)+A_P(0,0).
\]
All four references must use the same rendering path. Replacing only the clean
corner by a native reference reintroduces the offset from the existing paper.

\paragraph{Proposition 1 (algebraic joint-effect bound).}
For any norm, without differentiability assumptions,
\[
 \|d_{J,N}\|\ge\max\{0,\|I_N\|-\|d_{1,N}\|-\|d_{2,N}\|\}.
\]
If $\|I_R-I_N\|\le\xi$ and both native unilateral effects are at most
$\tau_1,\tau_2$, then
\[
 \|d_{J,N}\|\ge\max\{0,\|I_R\|-\xi-\tau_1-\tau_2\}.
\]
Proof: $d_{J,N}=I_N+d_{1,N}+d_{2,N}$ and the reverse triangle inequality.
RMS on a fixed $d$-dimensional action subset is a scaled Euclidean norm, so the
same bound holds for the experimental RMS. This is an action-effect bound,
not a task-failure or closed-loop safety certificate.

\paragraph{Proposition 2 (separable cross-path errors cancel).}
Let $E(u,v)=A_R(u,v)-A_N(u,v)$. Exactly,
\[
 I_R-I_N=E(\theta,\theta)-E(\theta,0)-E(0,\theta)+E(0,0).
\]
Consequently, if $E(u,v)=b+e_1(u)+e_2(v)$ at the four queried corners, then
$I_R=I_N$, even when the ordinary displacement error
$d_{J,R}-d_{J,N}=E(\theta,\theta)-E(0,0)$ is large.
This assumption is empirical, not a property guaranteed by dual rendering.
An error in nonlinear view fusion is not separable and need not cancel.

\paragraph{Proposition 3 (mixed-curvature representation).}
If $A_P$ is $C^2$ on the rectangle of texture interpolations, the fundamental
theorem of calculus twice gives
\[
 I_P=\int_0^1\!\int_0^1D_{12} A_P(s\theta,t\theta)
                    [\theta,\theta],dt\,ds.
\]
Thus, if the bilinear operator norm $\|D_{12}E\|\le L_{12}$ throughout it,
$\|I_R-I_N\|\le L_{12}\|\theta\|^2$.
This is an absolute second-order error bound. Since $I_N$ is itself generally
second-order, it does not guarantee small relative error or a useful gradient.
Actual PNG quantization, clipping, and finite-precision inference invalidate
unqualified $C^2$ claims; Propositions 1--2 still hold algebraically for their
four observed outputs. The smooth statement refers to an ideal continuous
rendering/policy extension.

\paragraph{Local spectral interpretation (not a new eigensolver).}
For an ideal smooth path and a fixed action direction $q$, define
$H_q=\tfrac12(D_{12}(q^\top A)(0,0)+D_{12}(q^\top A)(0,0)^\top)$.
Then $q^\top I(t v)=t^2v^\top H_qv+O(t^3)$ when third derivatives are bounded.
Let $J_1=D_1A(0,0)$, $J_2=D_2A(0,0)$, and $P$ project onto
$\ker J_1\cap\ker J_2$. The maximizer of $v^\top H_qv$ over unit vectors in
that subspace is a largest-eigenvalue vector of the restricted operator
$P H_q P$. The absolute-value version uses the largest-magnitude eigenvalue.
This does not solve $\max_v\|D_{12}A[v,v]\|$ globally; that vector objective
is quartic. With exact first-order unilateral nulling, both solo outputs and
the interaction become second-order, so finite-step constraints remain needed.

\paragraph{Candidate optimization.}
One possible source-only test is
\[
 \max_\theta\ \mathbb E_{x,u}\|I_R(\theta;x,u)\|^2,
 \quad C_v(\theta)=\mathbb E_{x,u}\|d_{v,R}(\theta;x,u)\|^2\le\tau_v^2,
 \quad T_\theta\in\mathcal T_\epsilon.
\]
The thresholds are fixed source-side and must be checked per state or via a
prespecified tail constraint if worst-state protection is claimed. An average
constraint alone does not give each-state protection. All original amplitude,
TV, rasterization/export, and source/target separation requirements remain.

For $g=\nabla\|I_R\|^2$ and $h_v=\nabla C_v$, an economical local subproblem is
\[
 \max_{\|\delta\|_2\le r}g^\top\delta
 \quad\text{subject to }h_v^\top\delta\le0,\ v=1,2.
\]
Writing $\mathcal C$ for this convex cone, its exact solution is
$\delta=r\Pi_{\mathcal C}(g)/\|\Pi_{\mathcal C}(g)\|$, or zero when the
projection vanishes. The projection has a two-dimensional nonnegative dual;
enumerating zero, one, or two active constraints solves the small problem.
This is standard gradient-constrained quadratic programming, closely related
to Gradient Episodic Memory, not an independent mathematical novelty claim.
If $C_v$ has $\beta_v$-Lipschitz gradient, a feasible current point satisfies
$C_v(\theta+\delta)\le C_v(\theta)+\beta_v r^2/2$.
In practice use a finite-step backtracking check with frozen margins; do not
infer protection from the sign of a local derivative alone. A zero/small
initialization has a vanishing interaction gradient, so spectral or nonzero
initialization needs to be separated from training gains in the comparisons.
